\chapter{Discussion}
\label{chp:discussion}
% Shortened on 19 Nov 2026 (previous version: scratchpad 8.discussion.before_cut.tex).

This chapter interprets the results against the question that motivates the study,
whether a simulation learned from the event log of CTB can support the assessment of
a different block allocation for truck deliveries. It discusses what the model can
and cannot support, what the case shows about discovery metrics and about
data-aware simulation in a physical process, which threats remain, how domain
knowledge entered the study, and what the terminal can take from it.

\section{Adequacy for the Intended Use}
\label{sec:disc_adequacy}

Three properties decide whether a simulation can support the comparison of a block
allocation rule with the current allocation. It must generate only visits that
trucks can perform, it must produce the waiting at a busy block by the mechanism that
the allocation changes, and it must respond to load as strongly as the terminal. The
final configuration meets the first two and the third only in part.

\paragraph{What the model supports.} Every simulated visit satisfies the checked
rules, and the waiting before a stop arises when all three lanes of a block are
taken. The model can therefore show how moving deliveries away from full blocks acts
in principle, which visits gain, and whether the waiting is reduced or only moved.
Its result for the allocation rule is a trade-off between how often and how badly
the corridor is exceeded, and between the trucks that gain and those that pay
(Section~\ref{sec:res_s1_contrasts}). A plausible explanation for why the trade-off
appears only under higher demand lies in the mechanism. With moderate congestion, a
moved delivery finds a free lane, and the trucks at the receiving block hardly notice
it. Under heavy congestion, the receiving blocks fill as well, and a wider choice
moves more deliveries into them.

\paragraph{What the model does not support.} It does not give the size of the effect
at the terminal. It produces less than half of the hours that the terminal counts as
too long, its stops slow down less than the recorded ones at the busiest blocks, and
how strongly it congests depends on the depth of its trees
(Section~\ref{sec:res_e4_depth}). The size of the simulated effect depends on exactly
these properties (Section~\ref{sec:res_s1_sensitivity}). The model also leaves out
what a move costs at the terminal, namely the driving to another block, a later
rehandling of the container and the constraints of the yard plan. These can make the
real effect smaller, while a stronger slowdown under load can make it larger. The
supported statement is therefore conditional. If the waiting at a full block arises
as in the model, if eligible blocks exist and if a move costs little, then moving
deliveries away from blocks that three trucks are already heading to shortens the
long visits and lowers the share of hours above the corridor. Because the model does
not reproduce the severe congestion that motivates the rule, even the direction of
the effect at CTB remains a hypothesis that a pilot would have to test. Only the
upper bound of the corridor is evaluated, because the shift of equipment from HO2 to
the waterside that the lower bound allows is not recorded
(Section~\ref{sec:prob_statement}).

\paragraph{Model error and drift.} Part of the shortfall comes from the change
between the periods. The stops of the test period took longer at every load, and more
blocks were closed (Sections~\ref{sec:res_e4_remaining} and~\ref{sec:obs_drift}). One
weakness, however, appears in both periods. The model produces about half of the hours
above the corridor even in-sample. The missing slowdown in the busiest hours is
therefore a limitation of the model, not of the test period, and it lies in the
congestion that the allocation rule is meant to relieve.

\paragraph{Congestion outside the terminal.} The highest observed peaks are not
driven by the load at the blocks. In the test period they occur at the usual number
of arrivals, and the additional time lies in the exit after the last stop
(Section~\ref{sec:res_e4_remaining}). In discussions with the terminal during the
development of the model, these peaks were attributed to traffic congestion outside
the terminal, which delays the trucks on their way out and which the terminal cannot
influence. Such episodes need not be simulated as individual events, but they recur,
and how the terminal copes with them is of particular operational interest. The
simulation cannot generate them, because they are rare and their cause is not
recorded in the event log. The three episodes of the test period contain only
11\,\% of the hours above the corridor, but 36\,\% of the minutes by which these hours
exceed it. They therefore explain a large part of the missing severity, but little of
the missing frequency, and they count against the corridor although the terminal
cannot influence them. A different allocation of deliveries cannot relieve congestion
outside the terminal, and the model cannot show how the operation, or a rule,
performs during such an episode.

\paragraph{Conditions of the result.} The result is conditional on the replayed
arrivals, case attributes and blocks. Without them, the cycle-time distance rises from
8.0 to 11.2 minutes, and a prospective use would need an arrival model with the daily
peaks and the block of each container from the yard plan. It also rests on the
extensions of the simulator, each of which the load response needs
(Section~\ref{sec:res_e4_loo}). Two details of the implementation bound it further.
The block is replayed only for the 75\,\% of the RMG stops that have an observed
counterpart, and the load that the service time sees is counted before a stop may
move to a free lane (Appendix~\ref{app:verification}).

\section{Discovery Metrics and Simulation Behaviour}
\label{sec:disc_metrics}

Good discovery metrics predicted a usable simulation only in part. First, what a net
permits is not what the simulator generates. ProSiT learns its routing from the
observed visits \cite{Vinci2026prosit}, so behaviour that a net permits but the log
never shows is generated rarely, and precision separated the nets mainly where they
are very permissive. Second, coverage of the frequent behaviour mattered throughout,
because a sequence that a net rejects is never generated. IMf 0.3 loses almost all
visits to two RMG blocks, which are long visits and the ones a block allocation
affects most. Third, no measure on the language shows how a net is built, and the
structure decided what the simulator does. IM 0.0 has parallel branches that silent
transitions can skip, and in the simulation 12 to 21\,\% of its multi-stop visits have
overlapping stops, because ProSiT learns each routing weight from the case attributes
and the history of the case, not from the marking of the net
(Appendix~\ref{app:replica}). The same holds for the stochastic labelled data Petri
nets on which ProSiT builds
\cite{mannhardt2023stochasticProcessesModelling,Vinci2023Inf.Activ.Prob}. Under the
final configuration, these impossible visits even bring IM 0.0 closest to the observed
turnaround time (Section~\ref{sec:res_e4_nets}), so a selection by the simulated
turnaround time alone would have chosen a net that no truck can follow. In HM 0.7,
every decision sits on a silent transition, and a data rule has to look ahead on the
net. Conformance metrics evaluate the control flow, not the resource mechanisms, so
the two strongest effects in the simulation, the self-count in the queue feature and
the count of trucks at a block, could only be found by validating the simulator
component by component.

\paragraph{What the discovered control flow contributes.} The control flow decides
which sequences the simulator can generate at all, whether impossible visits arise,
where a data rule has to act and at which decision points ProSiT learns its routing.
Once the resource mechanisms are extended, however, nets without parallel branches
give nearly the same turnaround time, and the net explains at most 1.5 of the
8.0 minutes of cycle-time distance (Section~\ref{sec:res_e4_nets}). The discovered
control flow is therefore a necessary basis of the simulation, but not the source of
its remaining error.

\paragraph{Simplicity and the analyst.} Simplicity was not used for the selection,
because in this log a small net obtains its size either by rejecting frequent
sequences or by allowing every sequence of stops (Section~\ref{sec:res_rules}). The
selected net has a cost for the analyst, however. With 34 silent transitions, HM 0.7
is harder to read and to change by hand than its language suggests, and an analyst
who adds an activity or changes a rule has to understand at which silent transitions
the decisions are taken. The stop guard is therefore computed from the reachability
graph and recomputed when the net changes, and the same language encoded as a state
machine generates the same visits (Appendix~\ref{app:replica}), a representation that
an analyst may find easier to read. For a model that analysts are meant to maintain,
the size of the net is a criterion of its own, but one to weigh only among nets that
already cover the frequent behaviour and respect the domain rules.

For model selection in simulation studies, this suggests an order. First the coverage
of the frequent behaviour, then the domain rules and precision, then a check of how
the net expresses its decisions, and finally a validation layer by layer in the
simulator itself. This agrees with calls for validation procedures that draw on
simulation methodology rather than on conformance measures alone
\cite{Sulis2026SimulationReview}, and with the critique of precision measures
\cite{Tax2018Imprecisions} and of log-similarity evaluation
\cite{Oezdemir2025RethinkingBPS}. The evidence comes from one process with a simple
control flow and one discovery grid. It shows that the metrics can fail to anticipate
the simulation, not how often they do.

\section{Data-Aware Simulation in a Physical Logistics Process}
\label{sec:disc_transfer}

\paragraph{What transfers.} Much of ProSiT transferred without change. The routing
trees split on the case attributes and reproduce the stop structure, and 98.7\,\% of
the simulated visits have their planned number of stops even without the guard
(Section~\ref{sec:res_baseline}). The waiting and execution trees use plausible
features such as the queue, the workload, the hour and reefer containers. The
white-box form was a precondition for the mechanism audit, because every finding of
E4 started from a tree that could be read.

\paragraph{What does not transfer out of the box.} The findings of E4 are of four
kinds. The first three are causes of the discrepancies named in
Section~\ref{sec:res_rq2}. The fourth cause named there, the change between the
periods, belongs to this data set, and in its place a limit appears that concerns
every setting, the validity under an intervention.
\begin{itemize}
\item \textbf{Implementation.} The queue feature counts the stop itself in training
but not in the simulation, and the count of stops in progress is too low for 14 to
15\,\% of the RMG stops. These defects can be corrected in the code and say little
about the method. Version 1.0.5 removes the self-count \cite{ProSiTRepository}, but
on its own this leaves the turnaround time without a response to load
(Section~\ref{sec:res_e4_loo}).
\item \textbf{Resource semantics.} The waiting is measured from a completely free
block, the capacity is checked when a stop becomes possible, a full block makes a
stop wait for the stop booked last, and the service time has no load feature. These
choices fit a resource that works through its bookings one after another, not a
block with three lanes. They match the known risks of business process simulation,
the neglected context and the working speed that depends on the workload
\cite{Aalst2015}, and the gap for queues and resources that Sulis et al.\ find
\cite{Sulis2026SimulationReview}. Queue mining asks these questions from the start
\cite{Senderovich2015QueueMining}.
\item \textbf{Data.} The log does not record the use of the cranes, their competition
with the waterside and the rail, or the waiting for a lane within the recorded stop.
\item \textbf{Validity under intervention.} The learned relations between the load of
a block and the waiting and service are associations. Whether they hold when the
allocation changes cannot be checked with these data.
\end{itemize}

\paragraph{What had to be added.} The simulation model keeps learned everything that
the log can express and adds what it cannot (Table~\ref{tab:sim_adaptations}). The
number of lanes, the domain rules and the decision point of the allocation come from
the terminal. The waiting target, the load feature and the exact count change how
the training data are computed, while the trees are still learned with ProSiT's own
functions, settings and random state. The values of the predictive parameters are
therefore learned, and only their meaning was chosen to match the physical process.
This is the split that data-driven simulation in manufacturing uses, where the
structure is discovered and buffers and capacities are added
\cite{LugaresiMatta2021,Grobis2025ProductionLogistics}, and extending a discovered
model where its assumptions do not fit is common in business process simulation
\cite{RWChapela2024Delays,runtimeIntegration2024}. With all extensions switched off,
the simulator reproduces ProSiT exactly, so the findings about ProSiT itself remain
visible.

\paragraph{What a case-centric log cannot show.} The cranes of a block also serve
the waterside and the rail, containers are restacked and trucks drive between areas,
and in the log all of this appears only as longer stops and approaches. An
object-centric event log with trucks, containers, blocks and cranes as objects
\cite{VanDerAalst2023OCPM} and object-centric simulation models
\cite{Knopp2023ObjectCentricSimulation} would be a next step. The most direct gain
would come from one more event per stop. If the log recorded when the crane starts to
handle a truck, the waiting for the crane would be an ordinary waiting time at a
resource that handles one truck after another, the case for which ProSiT was built.

\paragraph{Beyond the terminal.} For other processes in which several cases share a
capacity-limited resource, the findings suggest checking whether training and
simulation compute the workload in the same way, whether the waiting-time target
matches the resource, and whether the learned durations cover the load conditions
that the intervention creates. These are hypotheses from one terminal, one control
flow and one version of ProSiT, not evidence that every such process needs the same
extensions.

\section{Threats to Validity}
\label{sec:disc_threats}

The threats are grouped as construct, internal, external and conclusion validity
\cite{Wohlin2012Experimentation}. The design already answers some of them
(Section~\ref{sec:exp_threats}). This section states what remains.

\subsection{Construct Validity}
\paragraph{What a stop measures.} A recorded stop is time at the block, not service
time. It can contain waiting for the crane, and the log records a fourth truck at a
block for 5.6\,\% of the training stops, which may be a truck waiting for a lane or a
delayed completion record. The three lanes of the model are therefore an
approximation of the physical limit. Part of the waiting in the recorded stops stays
in the learned service time, which may be one reason why the simulated stops grow too
little at the busiest blocks. A truck that had to queue for a lane also loses the
learned remainder of its waiting, which may contribute to the shortfall of the
approach at a full block.

\paragraph{Cleaning and thresholds.} The cleaning removed 1.4\,\% of the visits as
recording errors, and the coverage threshold of E1 leaves about 1\,\% of the training
visits out of the behaviour that a net must accept. Both decisions are reproducible,
but not neutral, because they shape the long tail of the turnaround time that the
discovery sees and the simulation learns (Section~\ref{sec:cleaning_population}).

\paragraph{The KPI.} The corridor applies to the mean turnaround time of the trucks
that leave in an hour (Section~\ref{sec:prob_statement}), and only hours with at least
20 trucks count. Without this minimum, the shares of hours above 35 minutes change by
at most 2.4 percentage points, but the lowest weekly rank correlation of the test
period rises to 0.54, which the final configuration would miss. The approach mixes
driving and waiting, timestamps are floored to the minute, and the validity of the
model is evaluated for the turnaround time and its parts only, not for the use of
cranes or lanes.

\paragraph{The case attributes.} The case attributes come from the planned visit,
which is known at Gate In (Section~\ref{sec:entry_known_attributes}). If a schedule
were changed after Gate In, the planned number of stops would carry information from
the visit itself, and the stop guard would make the simulated number of stops look
better than a prospective run could.

\subsection{Internal Validity}
\paragraph{Retrospective hold-out.} Every selection uses the training log, and its
criteria were fixed before the runs they decide. The split itself was fixed after the
data had been explored, and the acceptance thresholds draw partly on the weekly
variation of the test period (Section~\ref{sec:exp_partitions}). The extensions rest
on mechanisms that the training data show, not on parameters fitted to the test
period, but the test error may still be somewhat optimistic. The error on the
congested test period mixes the error of the model with the change of the terminal,
and the in-sample run helps to distinguish the two without separating them cleanly.

\paragraph{Own code.} The extensions are switches in ProSiT's simulator, and with all
of them switched off the simulator reproduces ProSiT's log exactly
(Section~\ref{sec:sim_adaptations}). Conformance values and the soundness test come
from ProM. The measures of the generated logs and the control-flow replica are own
code.

\subsection{External Validity}
The data cover one terminal and eight and a half weeks in spring 2026. Seasonal
patterns, other terminals and other equipment are not covered, and the findings refer
to ProSiT~1.0.4. The discovery grid has 19 nets of four families, and other
algorithms such as the Split Miner \cite{RWSplitMiner2019} were not tested. The
mechanisms found here are candidates to check in other capacity-limited processes
(Section~\ref{sec:disc_transfer}), while the numbers are specific to the case.

\subsection{Conclusion Validity}
Each configuration was run with five seeds and each scenario of E5 with ten paired
seeds. The intervals describe the randomness of the simulator for one fitted model,
not the uncertainty of the fitted parameters or of the model structure. The
acceptance thresholds come from few weeks and two periods and are themselves
uncertain. The what-if analysis assumes that the learned associations between load,
waiting and service hold when the allocation changes, which the data cannot confirm
\cite{Vinci2026prosit}.

\section{Domain Knowledge and Stakeholder Review}
\label{sec:disc_stakeholders}

The terminal took part in the study in two ways. Its data warehouse and business
intelligence team and terminal development co-designed the use case in every phase,
from the data and the cleaning to the domain rules and the improvement measure. In
addition, the completed model and its results were reviewed by e-mail on 12 November
2026 by the project manager responsible for the cross-terminal KPI dashboard, who is
also the company-side co-supervisor of this thesis, so the review is not independent
of the study.

The review confirmed the domain rules, the three lanes per block and that all
sequences of the event log are real operational variants. It confirmed that the yard
position of a delivery is assigned at Gate In, near the berth of the target vessel if
it is known and in block T13 otherwise, without considering how many trucks are
already heading to a block. This is the decision that the improvement measure changes.
The review corrected the corridor, which applies to the mean turnaround time of the
trucks that leave in an hour, and pointed out that below the corridor the terminal can
move manual equipment from HO2 to the waterside. It evaluated the result of the
improvement measure as plausible, because moves within one or two blocks need no
rehandling and no long additional driving, and it named a completely free block
choice as an open question. It did not find the missing peak hours of the simulation
plausible, which agrees with the main limitation found in E4. These statements are an
expert assessment, not a measured cost. An assumption that the modeller and the
terminal put into the model cannot be confirmed by the same model. The records can
only show that they do not contradict it, as they do for rules 1 to 5.

\section{Implications for Practice}
\label{sec:disc_practice}

\paragraph{Where the time is lost.} Visits longer than 35 minutes are longer in every
part of the visit, before the first stop, at the stops and before Gate Out
(Table~\ref{tab:kpi_long_visits}), and the waiting before the first stop grows
sharply when three trucks are already at a block. This locates two levers, the load
of the blocks at the moment a truck arrives and the waiting before Gate Out, part of
which comes from congestion outside the terminal.

\paragraph{What a block allocation can do.} The model points to a lever that
today's assignment does not use, the trucks already heading to each block when a
delivery passes Gate In. A local rule acts at the moment at which the terminal
assigns the yard position and moves a minority of the deliveries by one or two block
numbers. A wider choice of blocks removes more hours above the corridor, but under
high demand it shifts waiting to the trucks at the receiving blocks, while the
narrowest rule is the more cautious candidate (Section~\ref{sec:res_s1_contrasts}).

\paragraph{What would be needed before acting.} A decision to change the allocation
would need data on the occupancy of the lanes and cranes, so that the slowdown in the
busiest hours can be learned, a check in the yard plan of how often a nearby block can
really take a given container, and a small field trial against which the simulated
effect could be checked. Until then, the simulation is a tool to explain where
waiting arises and how a rule would act on it, not a forecast of its effect.
